% Lösungen Einheit 4 von 4 – Parallele Ebenen (zusammenbau v0.1)
\einheitenkopf{Einheit 4 von 4 · Parallele Ebenen}

\erg{17}{a) identisch \quad b) $F\colon 3x + y - 2z = 5$ \quad c) $T\colon 3x + 2y + 3z = 21$ (mit $K$); $L$: $3 + 0 + 18 = 21$ ✓, $L$ liegt in $T$ (in $T$ einsetzen, nicht in $S$) \quad d) $(4 | 2 | 4) = 2 \cdot (2 | 1 | 2)$: parallel oder identisch; $(0 | 8 | 0)$ liegt in $L_1$, in $L_2$ ergibt sich $2 \cdot 8 = 16 \ne 14$ – die Ebenen sind echt parallel \quad e) $E_B\colon x - y + 6z = 6$; die Höhe $z = \frac{6 - x + y}{6}$ ist über $D$ am größten ($x = 0$, $y = 6$): $z = 2$ m (über $A$ nur 1 m) \quad f) Volumen $2 : 1$ heißt Höhen $2 : 1$, $M$ schneidet $BE$ nach einem Drittel von $B$ aus: $(1 | 1 | 0) + \frac{1}{3} \cdot (3 | 3 | 6) = (2 | 2 | 2)$; $M\colon x + 2y + 2z = 10$ (nicht die halbe Höhe)}
\erg{18}{a) I und IV scheiden aus (Normalenvektor nicht kollinear zu $(2 | 0 | 5)$); für gleiches $x$ wächst $2x + 5z$ mit der Höhe, also mittlere Etage $2x + 5z = 50$, obere $2x + 5z = 60$; II läge unter der unteren Etage}
\erg{19}{a) Paul setzt $L$ in $E$ statt in $T$ ein; $T\colon 2x + y - z = 4$, $L$: $4 + 1 - 1 = 4$ ✓, $L$ liegt in $T$ \quad b) Parallele Ebenen haben dieselben Richtungen; ein Vektor, der auf diesen Richtungen senkrecht steht, steht auf beiden Ebenen senkrecht. \quad c) $F\colon x + 3z = 27$; bei $x = 0$ liegt $F$ in $z = 9$ m Höhe, das Dach in 10 m – die Schicht liegt dort 1 m tiefer}
